\begin{workedexample}[Worked Example: Loaded Q for a channel filter]
A bandpass filter at $f_0 = 2\ \text{GHz}$ must pass a $20\ \text{MHz}$
channel. Its loaded quality factor is
\[ Q_L = \frac{f_0}{\text{BW}} = \frac{2\times10^{9}}{20\times10^{6}} = 100. \]
For low insertion loss the resonators' \emph{unloaded} $Q$ must far exceed this
($Q_u \gg Q_L$). Microstrip ($Q_u\!\sim\!200$) would be lossy here; a cavity or
BAW resonator ($Q_u$ in the thousands) is the right choice.
\end{workedexample}

\begin{keytakeaways}
\begin{itemize}[leftmargin=1.4em,itemsep=2pt]
\item Distributed (coupled-line/combline/interdigital) filters print in microstrip; no lumped parts.
\item Cavity/DR give $Q_u$ in the thousands (low loss, sharp); SAW/BAW pack steep filters into a chip.
\item Insertion loss falls as $Q_u/Q_L$ rises; duplexers/diplexers share one antenna across TX/RX or bands.
\end{itemize}
\end{keytakeaways}

\begin{exercises}
\item Loaded $Q$ for a $1\%$ fractional-bandwidth filter?
\item Which filter technology sits in a smartphone's RF front end for steep, tiny filters?
\item Why does a base-station duplexer use cavity resonators rather than microstrip?
\end{exercises}

\furtherreading{Hong~\cite{hong}; Pozar~\cite{pozar} (ch.~8).}
